% GENERATED FILE. Edit canonical lessons, then run npm run generate:books.
% canonical-source-hash: fnv1a64:86e589d85d50e5af
% canonical-lessons: AR-C477-qirsh, AR-C477-saltaun, AR-C477-mahar, AR-C477-naqi, AR-C477-madiq

\chapter{Under the sea}
\label{ch:ar-under-the-sea}
\hlchaptermodality{477}

\begin{chapteropening}
\textbf{By the end of this chapter:} \emph{I can say a shark, a crab, oysters, pure and a strait.}

Complete the last lesson of chapter 477: I can say a shark, a crab, oysters, pure and a strait.
\end{chapteropening}

\section[\texorpdfstring{\ar{قرش} (qirsh)}{قرش (qirsh)}]{\ar{قرش} (qirsh) — a shark}

\label{lesson:AR-C477-qirsh}



\begin{quote}
Before the new one: say the Arabic for a claw, then the Arabic for to prey on.
\end{quote}

\subsection*{You'll want to know: \ar{قرش}}
\textbf{\ar{قرش}} — \emph{qirsh} — \textquotedblleft{}a shark\textquotedblright{}.

\begin{quote}
\ar{رأى} \ar{الصيادون} \ar{قرشا} \ar{بجانب} \ar{القارب}.
\end{quote}

\begin{quote}
\emph{raʾā aṣ-ṣayyādūn qirshan bi-jānib al-qārib.}
\end{quote}

\begin{quote}
The fishermen saw a shark beside the boat.
\end{quote}

\emph{qirsh} is also an old coin, a piastre.

\begin{cousinweb}
\textbf{\ar{قرش}} (\emph{qirsh}) is usually traced to the native Arabic root \emph{q-r-sh}, \textquotedblleft{}to gather, to scrape together\textquotedblright{}.
\end{cousinweb}

\begin{grammarlens}[title={What you've built}]
One more word for getting around.
\end{grammarlens}

\subsection*{Your turn}
\begin{itemize}
\raggedright
  \item \emph{Say it:} \emph{qirsh}
  \item \emph{Say it:} \emph{qirsh}, once more
  \item \emph{Say it:} \emph{raʾā aṣ-ṣayyādūn qirshan bi-jānib al-qārib.} — The fishermen saw a shark beside the boat.
  \item \emph{Recall:} say the Arabic for a claw, then the Arabic for to prey on, then say \emph{qirsh} again
\end{itemize}

\begin{tcolorbox}[breakable,colback=teal!4,colframe=teal!35!black,title={Before you move on}]
What does \ar{قرش} mean? (\textquotedblleft{}A shark\textquotedblright{}.) Say it once more.
\end{tcolorbox}
\section[\texorpdfstring{\ar{سلطعون} (salṭaʿūn)}{سلطعون (salṭaʿūn)}]{\ar{سلطعون} (salṭaʿūn) — a crab}

\label{lesson:AR-C477-saltaun}



\begin{quote}
Before the new one: say the Arabic for to prey on, then the Arabic for a shark.
\end{quote}

\subsection*{You'll want to know: \ar{سلطعون}}
\textbf{\ar{سلطعون}} — \emph{salṭaʿūn} — \textquotedblleft{}a crab\textquotedblright{}.

\begin{quote}
\ar{أمسك} \ar{أخي} \ar{سلطعونا} \ar{على} \ar{الصخرة}.
\end{quote}

\begin{quote}
\emph{amsaka akhī salṭaʿūnan ʿalā aṣ-ṣakhra.}
\end{quote}

\begin{quote}
My brother caught a crab on the rock.
\end{quote}

\begin{cousinweb}
\textbf{\ar{سلطعون}} (\emph{salṭaʿūn}) names the crab; the origin is not settled.
\end{cousinweb}

\begin{grammarlens}[title={What you've built}]
One more word for getting around.
\end{grammarlens}

\subsection*{Your turn}
\begin{itemize}
\raggedright
  \item \emph{Say it:} \emph{salṭaʿūn}
  \item \emph{Say it:} \emph{salṭaʿūn}, once more
  \item \emph{Say it:} \emph{amsaka akhī salṭaʿūnan ʿalā aṣ-ṣakhra.} — My brother caught a crab on the rock.
  \item \emph{Recall:} say the Arabic for to prey on, then the Arabic for a shark, then say \emph{salṭaʿūn} again
\end{itemize}

\begin{tcolorbox}[breakable,colback=teal!4,colframe=teal!35!black,title={Before you move on}]
What does \ar{سلطعون} mean? (\textquotedblleft{}A crab\textquotedblright{}.) Say it once more.
\end{tcolorbox}
\section[\texorpdfstring{\ar{محار} (maḥār)}{محار (maḥār)}]{\ar{محار} (maḥār) — oysters; shellfish}

\label{lesson:AR-C477-mahar}



\begin{quote}
Before the new one: say the Arabic for a shark, then the Arabic for a crab.
\end{quote}

\subsection*{You'll want to know: \ar{محار}}
\textbf{\ar{محار}} — \emph{maḥār} — \textquotedblleft{}oysters; shellfish\textquotedblright{}.

\begin{quote}
\ar{أكلنا} \ar{المحار} \ar{في} \ar{مطعم} \ar{على} \ar{البحر}.
\end{quote}

\begin{quote}
\emph{akalnā al-maḥār fī maṭʿam ʿalā al-baḥr.}
\end{quote}

\begin{quote}
We ate oysters at a restaurant by the sea.
\end{quote}

\begin{cousinweb}
\textbf{\ar{محار}} (\emph{maḥār}) names oysters and other shellfish; the origin is not settled.
\end{cousinweb}

\begin{grammarlens}[title={What you've built}]
One more word for getting around.
\end{grammarlens}

\subsection*{Your turn}
\begin{itemize}
\raggedright
  \item \emph{Say it:} \emph{maḥār}
  \item \emph{Say it:} \emph{maḥār}, once more
  \item \emph{Say it:} \emph{akalnā al-maḥār fī maṭʿam ʿalā al-baḥr.} — We ate oysters at a restaurant by the sea.
  \item \emph{Recall:} say the Arabic for a shark, then the Arabic for a crab, then say \emph{maḥār} again
\end{itemize}

\begin{tcolorbox}[breakable,colback=teal!4,colframe=teal!35!black,title={Before you move on}]
What does \ar{محار} mean? (\textquotedblleft{}Oysters; shellfish\textquotedblright{}.) Say it once more.
\end{tcolorbox}
\section[\texorpdfstring{\ar{نقي} (naqī)}{نقي (naqī)}]{\ar{نقي} (naqī) — pure; clean}

\label{lesson:AR-C477-naqi}



\begin{quote}
Before the new one: say the Arabic for a crab, then the Arabic for oysters.
\end{quote}

\subsection*{You'll want to know: \ar{نقي}}
\textbf{\ar{نقي}} — \emph{naqī} — \textquotedblleft{}pure; clean\textquotedblright{}.

\begin{quote}
\ar{الماء} \ar{هنا} \ar{نقي} \ar{جدا}.
\end{quote}

\begin{quote}
\emph{al-māʾ hunā naqī jiddan.}
\end{quote}

\begin{quote}
The water here is very pure.
\end{quote}

\begin{cousinweb}
\textbf{\ar{نقي}} (\emph{naqī}) is built on the native Arabic root \emph{n-q-y}, \textquotedblleft{}to be clean\textquotedblright{} (Lane).
\end{cousinweb}

\begin{grammarlens}[title={What you've built}]
One more word for getting around.
\end{grammarlens}

\subsection*{Your turn}
\begin{itemize}
\raggedright
  \item \emph{Say it:} \emph{naqī}
  \item \emph{Say it:} \emph{naqī}, once more
  \item \emph{Say it:} \emph{al-māʾ hunā naqī jiddan.} — The water here is very pure.
  \item \emph{Recall:} say the Arabic for a crab, then the Arabic for oysters, then say \emph{naqī} again
\end{itemize}

\begin{tcolorbox}[breakable,colback=teal!4,colframe=teal!35!black,title={Before you move on}]
What does \ar{نقي} mean? (\textquotedblleft{}Pure; clean\textquotedblright{}.) Say it once more.
\end{tcolorbox}
\section[\texorpdfstring{\ar{مضيق} (maḍīq)}{مضيق (maḍīq)}]{\ar{مضيق} (maḍīq) — a strait}

\label{lesson:AR-C477-madiq}



\begin{quote}
Before the new one: say the Arabic for oysters, then the Arabic for pure.
\end{quote}

\subsection*{You'll want to know: \ar{مضيق}}
\textbf{\ar{مضيق}} — \emph{maḍīq} — \textquotedblleft{}a strait\textquotedblright{}.

\begin{quote}
\ar{عبرت} \ar{السفينة} \ar{المضيق} \ar{في} \ar{الليل}.
\end{quote}

\begin{quote}
\emph{ʿabarat as-safīna al-maḍīq fī al-layl.}
\end{quote}

\begin{quote}
The ship crossed the strait at night.
\end{quote}

\begin{cousinweb}
\textbf{\ar{مضيق}} (\emph{maḍīq}) is built on the native Arabic root \emph{ḍ-y-q}, \textquotedblleft{}to be narrow\textquotedblright{} (Lane). It is a narrow place, beside \emph{ḍayyiq}, \textquotedblleft{}narrow\textquotedblright{}.
\end{cousinweb}

\begin{grammarlens}[title={What you've built}]
One more word for getting around.
\end{grammarlens}

\subsection*{Your turn}
\begin{itemize}
\raggedright
  \item \emph{Say it:} \emph{maḍīq}
  \item \emph{Say it:} \emph{maḍīq}, once more
  \item \emph{Say it:} \emph{ʿabarat as-safīna al-maḍīq fī al-layl.} — The ship crossed the strait at night.
  \item \emph{Recall:} say the Arabic for oysters, then the Arabic for pure, then say \emph{maḍīq} again
\end{itemize}

\begin{tcolorbox}[breakable,colback=teal!4,colframe=teal!35!black,title={Before you move on}]
What does \ar{مضيق} mean? (\textquotedblleft{}A strait\textquotedblright{}.) Say it once more.
\end{tcolorbox}
